% ============================================================
% born_rule_from_paths.tex
% Section: The Born Rule as a Path-Counting Identity
% Corresponds to: src/experiments/exp_06_path_counting.py
%                 notes/born_rule_from_path_counting.md
% ============================================================

\subsection{The Born Rule as a Path-Counting Identity}
\label{subsec:born_rule_paths}

The Born rule --- the statement that the probability of finding
a session at node $\mathbf{x}$ is $|\psi(\mathbf{x})|^2$ --- is
a postulate in standard quantum mechanics.
In the $\unity$ framework it is a counting identity, with the
quadratic form forced by Gleason's theorem rather than chosen.
The path-counting derivation below gives the operational realisation
on the discrete lattice; we first record the uniqueness statement
that fixes the form to $|\psi|^2$ rather than $|\psi|^p$ for some
other $p$.

\subsubsection*{Uniqueness via Gleason}
\label{subsubsec:born_gleason}

The path-counting identity below shows $\rho = |\psi_R|^2 + |\psi_L|^2$
is \emph{consistent} with $\mathcal{A}=1$ on the bipartite tick rule.
That is sufficiency.  The remaining technical question is necessity:
why this measure and not, e.g., $|\psi|^p$ for $p \ne 2$?

The answer is Gleason's theorem~\cite{gleason1957}.  The per-session
state space $\mathcal{H}_\text{session} = \ell^2(\Tdiamond,\mathbb{C}^2)$
is a separable Hilbert space; for any practical lattice
$\dim \mathcal{H}_\text{session} \gg 3$.  The bipartite tick rule is
unitary, preserving the inner product
$\langle\psi, \phi\rangle = \sum_x (\psi_R^*\phi_R + \psi_L^*\phi_L)$.

A \emph{frame function} on the projection lattice
$\mathcal{P}(\mathcal{H}_\text{session})$ is a function
$\mu : \mathcal{P} \to [0,1]$ with $\mu(I) = 1$ and
$\sigma$-additivity on countable orthogonal families.  $\mathcal{A}=1$
gives normalisation; the parallelogram identity for orthogonal ranges
of unitary projections gives $\sigma$-additivity.  Both hypotheses are
satisfied --- and they are exactly the structural intuition that
motivated $\mathcal{A}=1$ in the first place: probability conservation
imposed at every tick on the simplest geometry capable of supporting
it produces a quadratic measure under unitary evolution.

Gleason's theorem then gives uniquely $\mu(P) = \mathrm{tr}(\rho P)$
for a unique density operator $\rho$, and for pure $\psi$ this is
$\|P\psi\|^2$.  Specialising to the projection onto site $\mathbf{x}$
summed over spinor components, $\mu(P_\mathbf{x}) =
|\psi_R(\mathbf{x})|^2 + |\psi_L(\mathbf{x})|^2$.  The Born rule is
forced.

The non-trivial point of Gleason is that the same is \emph{not} true
for non-orthogonal projection families: $\mu(P_1) + \mu(P_2) \ne
\mu(P_1 + P_2)$ when $P_1, P_2$ have non-orthogonal ranges.  This is
the source of interference in standard QM, inherited by the
bipartite-tick framework as a structural consequence rather than a
postulate.

Three scope caveats:

\begin{enumerate}
\item Multi-session events (\texttt{exp\_20} arm B beam splitter,
\texttt{exp\_19c} drain) move between Hilbert spaces of different
session count $N$.  The Born rule holds on each fixed-$N$ projection
lattice; transitions between are creation/annihilation operations
governed by the operation algebra (notes/lattice\_operation\_algebra.md,
follow-on paper~\#13).
\item The bipartite RGB/CMY parity gives a $\mathbb{Z}_2$ grading on
$\mathcal{H}_\text{session}$ but the projection lattice is unchanged;
Gleason applies as stated.
\item The path-counting derivation that follows is the
\emph{operational realisation} of the Gleason-uniqueness measure on
the discrete bipartite lattice --- it tells you how to compute
$|\psi|^2$ from Feynman-style path sums; Gleason tells you why the
result must be quadratic in $\psi$.
\end{enumerate}

\paragraph{Paths on the lattice.}
After $N$ ticks, a massless session ($\omega = 0$, $p_\mathrm{hop} = 1$)
on a flat lattice has taken one step along each of the six basis vectors
$\{\mathbf{V}_1, \mathbf{V}_2, \mathbf{V}_3, -\mathbf{V}_1, -\mathbf{V}_2,
-\mathbf{V}_3\}$ at each tick.
The total number of $N$-step paths from the origin is $6^N$.

A path reaches node $(a, b, c)$ if, for some non-negative integers
$n_i$ (steps along $\mathbf{V}_i$) and $m_i$ (steps along $-\mathbf{V}_i$)
with $\sum_i(n_i + m_i) = N$:
\begin{equation}
    \sum_{i=1}^{3} (n_i - m_i)\,\mathbf{V}_i = (a,b,c).
    \label{eq:path_constraint}
\end{equation}
Setting $p_i = n_i - m_i$ and solving with $\mathbf{V}_1 = (1,1,1)$,
$\mathbf{V}_2 = (1,-1,-1)$, $\mathbf{V}_3 = (-1,1,-1)$:
\begin{equation}
    p_1 = \tfrac{a+b}{2}, \quad
    p_2 = \tfrac{a-c}{2}, \quad
    p_3 = \tfrac{b-c}{2}.
    \label{eq:net_steps}
\end{equation}
These are integers when $a, b, c$ all have the same parity --- the
bipartite reachability condition.
Given the net steps $p_i$, the number of distinct $N$-step paths to
$(a,b,c)$ is
\begin{equation}
    \pathcount{N}{a}{b}{c}
    = \sum_{\substack{s_1+s_2+s_3=N \\ s_i \geq |p_i|,\;
           s_i \equiv p_i \!\pmod{2}}}
      \frac{N!}{\,n_1!\,m_1!\,n_2!\,m_2!\,n_3!\,m_3!\,},
    \label{eq:path_count}
\end{equation}
where $n_i = (s_i + p_i)/2$ and $m_i = (s_i - p_i)/2$.
This is a pure combinatorial integer for every $(N, a, b, c)$.

\paragraph{A=1 is the Born normalisation.}
Summing the path count over all reachable nodes exhausts every $N$-step path:
\begin{equation}
    \sum_{(a,b,c) \in \causalcone{N}} \pathcount{N}{a}{b}{c} = 6^N.
    \label{eq:path_partition}
\end{equation}
Dividing by $6^N$:
\begin{equation}
    \sum_{(a,b,c) \in \causalcone{N}} \frac{\pathcount{N}{a}{b}{c}}{6^N} = 1.
    \label{eq:born_normalisation}
\end{equation}
This is the $\unity$ constraint expressed combinatorially.
For a massless session on a flat lattice (no potential, $\omega = 0$), the
amplitude at $(a,b,c)$ after $N$ ticks is real and positive, and
\begin{equation}
    |\psi(a,b,c,N)|^2 = \frac{\pathcount{N}{a}{b}{c}}{6^N}.
    \label{eq:born_identity}
\end{equation}
The Born rule is not a measurement postulate applied to $\psi$.
It is the fraction of all $N$-step paths that reach $(a,b,c)$, normalised
by $6^N$, which itself is forced by $\mathcal{A} = 1$.

\paragraph{Continuum limit: the Gaussian probability.}
For large $N$ with $(a,b,c) = N(\alpha,\beta,\gamma)$ fixed, the sum in
Eq.~\eqref{eq:path_count} concentrates around $s_i \approx N/3$ and
the central limit theorem gives
\begin{equation}
    \frac{\pathcount{N}{N\alpha,N\beta,N\gamma}}{6^N}
    \;\longrightarrow\;
    (2\pi N)^{-3/2}\,|\det C|^{-1/2}\,
    \exp\!\left(-\frac{N}{2}\,(\alpha,\beta,\gamma)\,C^{-1}
    \begin{pmatrix}\alpha\\\beta\\\gamma\end{pmatrix}\right),
    \label{eq:gaussian_born}
\end{equation}
where $C = \tfrac{1}{6}\sum_{\text{all }6} \mathbf{v}\mathbf{v}^T$ is the
covariance of a single step.
This is the spreading Gaussian wavefunction of a free massless particle ---
the Born probability --- derived from path counting alone.

\paragraph{Massive sessions.}
For $\omega \neq 0$ each path acquires a phase $\prod_\text{steps}
e^{i\delta\phi/2}$ that depends on the local potential encountered.
The amplitude $\psi(a,b,c,N)$ is then a complex sum over phase-weighted
paths: the lattice realisation of the Feynman path integral~\cite{feynman1948}.
The Born rule $|\psi|^2$ receives both the path count (number of contributing
histories) and the interference between them (phase correlations).
In the long-wavelength limit $k \ll \pi/a$ the phases vary slowly and
the path-count Gaussian~\eqref{eq:gaussian_born} dominates; the standard
Born rule is recovered as a continuum approximation valid for $N \gg 1$.

\paragraph{Discrete corrections.}
For small $N$, the integer path count $\pathcount{N}{a}{b}{c}$ deviates from
the Gaussian prediction~\eqref{eq:gaussian_born}.
These are exact, computable corrections to the Born probability at short
distances.
Experiment~\texttt{exp\_06} confirms that such corrections exist and are
numerically accessible: the discrete fringe structure at small $N$ is
non-Gaussian and matches the path-count integer formula.
If the lattice spacing equals the Planck length, the corrections become
significant at $r \sim \ell_P$ and constitute a falsifiable departure from
standard quantum mechanics in the deep ultraviolet.

\paragraph{Summary.}
The Born rule is not a postulate about measurement outcomes.
It is the statement that probability IS the path fraction, and the path
fraction sums to one because every path must end somewhere.
$\mathcal{A} = 1$ is the combinatorial identity $\sum P(N,\mathbf{x}) = 6^N$,
rewritten as a physical conservation law.
The measurement axiom, usually the most conceptually unsatisfying postulate
of quantum mechanics, is here a tautology.
